\begin{proof}[Proof of \cref{obj:lem:sparse-likelihood}]
\begin{enumerate}
\item By \cref{obj:def:sparse-family}, the normalized atom probabilities are
\[
q_{ay,x}=\frac{\widetilde q_{ay,x}}{S(\mathbf z)}.
\]
Summing the four unnormalized masses in each cell gives
\[
\sum_{a,y\in\{0,1\}}\widetilde q_{ay,x}
=\frac{\epsilon(1-\epsilon)(1+z_x)+(1-\epsilon)^2+\epsilon^2z_x}{d}
=\frac{1-\epsilon+\epsilon z_x}{d}.
\]
Consequently,
\[
S(\mathbf z)
=1-\epsilon+\frac{\epsilon}{d}\sum_{x\in[d]}z_x
\ge 1-\epsilon\ge\frac12>0.
\]
Since \(d>0\) and \(z_x\ge0\), we also have
\[
1-\epsilon+\epsilon z_x>0.
\]
Summing the normalized atom probabilities therefore yields
\[
p_x=\frac{1-\epsilon+\epsilon z_x}{dS(\mathbf z)},
\qquad
\sum_{y\in\{0,1\}}q_{1y,x}
=\frac{\epsilon(1-\epsilon)(1+z_x)}{dS(\mathbf z)}.
\]
Taking their ratio, as in \cref{obj:def:observed-class}, gives
\[
\pi_x
=\frac{\epsilon(1-\epsilon)(1+z_x)}
{1-\epsilon+\epsilon z_x}.
\]
The overlap inequalities follow by clearing the positive denominator:
\[
\pi_x-\epsilon
=\frac{\epsilon(1-2\epsilon)z_x}{1-\epsilon+\epsilon z_x}\ge0,
\qquad
1-\epsilon-\pi_x
=\frac{(1-\epsilon)(1-2\epsilon)}
{1-\epsilon+\epsilon z_x}\ge0.
\]
Here \(1-2\epsilon\ge0\), \(z_x\ge0\), and \(1-\epsilon>0\).
Thus the sparse law satisfies \cref{obj:ass:shrinking-overlap} and belongs to the observed class of \cref{obj:def:observed-class}.
% lean: CausalSmith.Stat.OptvalueVanishingoverlapRate.sparse_likelihood

\item The preceding calculations establish the asserted cell probability and propensity. For the conditional outcome means, the arm masses and outcome-one atom probabilities are
\[
\begin{aligned}
\sum_{y\in\{0,1\}}q_{0y,x}
&=\frac{(1-\epsilon)^2+\epsilon^2z_x}{dS(\mathbf z)},
&
q_{01,x}
&=\frac{(1-\epsilon)^2+\epsilon^2z_x}{2dS(\mathbf z)},\\
\sum_{y\in\{0,1\}}q_{1y,x}
&=\frac{\epsilon(1-\epsilon)(1+z_x)}{dS(\mathbf z)},
&
q_{11,x}
&=\frac{\epsilon(1-\epsilon)z_x}{dS(\mathbf z)}.
\end{aligned}
\]
Both arm masses are strictly positive, because
\[
(1-\epsilon)^2+\epsilon^2z_x>0,
\qquad
\epsilon(1-\epsilon)(1+z_x)>0.
\]
Dividing the outcome-one probabilities by their respective arm masses, using \cref{obj:def:observed-class}, gives
\[
\mu_{0x}=\frac12,
\qquad
\mu_{1x}=\frac{z_x}{1+z_x}.
\]
These divisions remain valid at \(z_x=0\).
% lean: CausalSmith.Stat.OptvalueVanishingoverlapRate.sparse_likelihood

\item To evaluate the observed value, split each cell according to whether \(z_x\le1\) or \(z_x\ge1\). Since \(1+z_x>0\),
\[
z_x\le1\ \Longrightarrow\ \frac{z_x}{1+z_x}\le\frac12,
\qquad
z_x\ge1\ \Longrightarrow\ \frac{z_x}{1+z_x}\ge\frac12.
\]
In the first case \(\max\{z_x-1,0\}=0\), and in the second it equals \(z_x-1\). Hence, in both cases,
\[
\max\left\{\frac12,\frac{z_x}{1+z_x}\right\}
=\frac12+\frac{\max\{z_x-1,0\}}{2(1+z_x)}.
\]
Using the definition of \(\phi_\epsilon\) in the statement, the cell contribution becomes
\[
p_x\max\{\mu_{0x},\mu_{1x}\}
=\frac{p_x}{2}
+\frac{\phi_\epsilon(z_x)}{2dS(\mathbf z)}.
\]
Moreover, the normalizer identity from the first step gives
\[
\sum_{x\in[d]}p_x
=\frac{\sum_{x\in[d]}(1-\epsilon+\epsilon z_x)}
{dS(\mathbf z)}
=\frac{dS(\mathbf z)}{dS(\mathbf z)}
=1.
\]
Summing the cell contributions according to \cref{obj:def:observed-value} therefore yields
\[
\Psi(\mathbb P_{\mathbf z}^{\mathrm{sp}})
=\frac12+\frac{1}{2dS(\mathbf z)}
\sum_{x\in[d]}\phi_\epsilon(z_x).
\]
% lean: CausalSmith.Stat.OptvalueVanishingoverlapRate.sparse_likelihood

\item Fix a cell \(x\). For \(z\in[0,M]\) and \(a,y\in\{0,1\}\), define the Poisson means and their reference values by
\[
\lambda_{ay,x}(z):=Bn\widetilde q_{ay,x}(z),
\qquad
\lambda_{\star,ay,x}:=\lambda_{ay,x}(z_\star).
\]
Explicitly,
\[
\begin{aligned}
\lambda_{11,x}(z)&=\frac{Bn}{d}\epsilon(1-\epsilon)z,
&
\lambda_{10,x}(z)&=\frac{Bn}{d}\epsilon(1-\epsilon),\\
\lambda_{00,x}(z)&=\lambda_{01,x}(z)
=\frac{Bn}{2d}\bigl((1-\epsilon)^2+\epsilon^2z\bigr).
\end{aligned}
\]
These means are nonnegative throughout \([0,M]\). Since \(B>2\), \(n\ge1\), \(d\ge2\), and \(z_\star=M/2\ge1\), all four reference means are strictly positive. Powers at zero are interpreted with \(0^0=1\).

For each coordinate and every \(\tau\in\mathbb R\), the Poisson exponential series gives
\[
\begin{aligned}
E_{z_\star}\!\left[\tau^{N_{ay,x}}\right]
&=\exp(-\lambda_{\star,ay,x})
\sum_{\rho=0}^{\infty}
\frac{(\lambda_{\star,ay,x}\tau)^\rho}{\rho!}\\
&=\exp\!\left\{\lambda_{\star,ay,x}(\tau-1)\right\}.
\end{aligned}
\]
Here \(\rho\) ranges over the nonnegative integers; the series converges absolutely, as is seen by replacing \(\tau\) with \(|\tau|\). Combining the two likelihood factors for this coordinate and applying the identity with
\(\tau=\lambda_{ay,x}(z)\lambda_{ay,x}(w)/\lambda_{\star,ay,x}^2\) gives
\[
\begin{aligned}
&E_{z_\star}\!\left[
\left\{\exp(\lambda_{\star,ay,x}-\lambda_{ay,x}(z))
\left(\frac{\lambda_{ay,x}(z)}{\lambda_{\star,ay,x}}\right)^{N_{ay,x}}\right\}
\left\{\exp(\lambda_{\star,ay,x}-\lambda_{ay,x}(w))
\left(\frac{\lambda_{ay,x}(w)}{\lambda_{\star,ay,x}}\right)^{N_{ay,x}}\right\}
\right]\\
&\quad=
\exp\!\left\{
2\lambda_{\star,ay,x}-\lambda_{ay,x}(z)-\lambda_{ay,x}(w)
+\lambda_{\star,ay,x}
\left(\frac{\lambda_{ay,x}(z)\lambda_{ay,x}(w)}
{\lambda_{\star,ay,x}^2}-1\right)
\right\}\\
&\quad=
\exp\!\left\{
\frac{(\lambda_{ay,x}(z)-\lambda_{\star,ay,x})
(\lambda_{ay,x}(w)-\lambda_{\star,ay,x})}
{\lambda_{\star,ay,x}}
\right\}.
\end{aligned}
\]
% lean: CausalSmith.Stat.OptvalueVanishingoverlapRate.poisson_likelihood_gram_scalar

Independence of the four reference counts now implies
\[
E_{z_\star}[L_zL_w]
=\exp\!\left\{
\sum_{a,y\in\{0,1\}}
\frac{(\lambda_{ay,x}(z)-\lambda_{\star,ay,x})
(\lambda_{ay,x}(w)-\lambda_{\star,ay,x})}
{\lambda_{\star,ay,x}}
\right\}.
\]
The treated-success contribution is
\[
\frac{(\lambda_{11,x}(z)-\lambda_{\star,11,x})
(\lambda_{11,x}(w)-\lambda_{\star,11,x})}
{\lambda_{\star,11,x}}
=
\frac{Bn}{d}\frac{\epsilon(1-\epsilon)}{z_\star}
(z-z_\star)(w-z_\star).
\]
The treated-failure contribution equals zero because its mean is constant. The two control contributions sum to
\[
\sum_{y\in\{0,1\}}
\frac{(\lambda_{0y,x}(z)-\lambda_{\star,0y,x})
(\lambda_{0y,x}(w)-\lambda_{\star,0y,x})}
{\lambda_{\star,0y,x}}
=
\frac{Bn}{d}
\frac{\epsilon^4}{(1-\epsilon)^2+\epsilon^2z_\star}
(z-z_\star)(w-z_\star).
\]
Thus their total is \(\alpha(z-z_\star)(w-z_\star)\), with \(\alpha\) exactly as in the statement, proving the likelihood identity.
% lean: CausalSmith.Stat.OptvalueVanishingoverlapRate.sparse_likelihood

\item Finally, \(M>0\) and
\[
(1-\epsilon)^2+\epsilon^2M/2>0.
\]
Clearing these positive denominators, the required inner bound follows from
\[
\begin{aligned}
&2\epsilon^2\bigl((1-\epsilon)^2+\epsilon^2M/2\bigr)
-M\epsilon^4\\
&\qquad=2\epsilon^2(1-\epsilon)^2\ge0.
\end{aligned}
\]
Indeed, this is precisely the numerator obtained by subtracting the left-hand side from the right-hand side of
\[
\frac{\epsilon(1-\epsilon)}{M/2}
+\frac{\epsilon^4}{(1-\epsilon)^2+\epsilon^2M/2}
\le\frac{2\epsilon}{M}
\]
and multiplying by \(M\bigl((1-\epsilon)^2+\epsilon^2M/2\bigr)\).
Multiplication by the nonnegative factor \(Bn/d\), together with \(z_\star=M/2\), gives
\[
\alpha
\le\frac{Bn}{d}\frac{2\epsilon}{M}
=\frac{2Bn\epsilon}{dM}.
\]
% lean: CausalSmith.Stat.OptvalueVanishingoverlapRate.sparseGramCoefficient_bound
\end{enumerate}
\end{proof}
